\section{An exact position--value coloring formulation}
\label{sec:value-color}

This section gives an equivalent finite constraint model, not a decision
of any new order. In one view let each line $a$ be a bijection
$f_a:P\to Q$. For distinct $a,b$ the permutation
$\pi=f_b^{-1}f_a$ is fixed-point-free. A permutation has only even
cycles if and only if its points admit a binary coloring that toggles
along every edge: alternating a coloring around a cycle is consistent
exactly when its length is even.

\begin{theorem}[Position--value channel]
\label{thm:value-color-channel}
Introduce colors $e_i,h_v\in\{0,1\}$ for positions $i\in P$ and
values $v\in Q$. The equations
\begin{equation}
 e_i=h_{f_a(i)},\qquad e_i=1-h_{f_b(i)}\quad(i\in P)
 \label{eq:value-color-channel}
\end{equation}
are satisfiable if and only if $\pi$ has only even cycles. Every
solution has exactly $n/2$ ones in each color family. One may
add $e_{i_0}=0$ at any chosen position without changing satisfiability.
\end{theorem}

\begin{proof}
For $j=\pi(i)$, equality $f_a(i)=f_b(j)$ makes
\eqref{eq:value-color-channel} imply $e_j=1-e_i$.
Conversely, from a toggling coloring set $h_v=e_{f_a^{-1}(v)}$.
The first equation is immediate. The second follows because
$f_a^{-1}f_b=\pi^{-1}$ and toggling also holds on inverse edges.
Since the lines are bijections, summing the two equations gives
$\sum_i e_i=\sum_v h_v=n-\sum_v h_v=n/2$.
Complementing all $e$ and $h$ simultaneously preserves the
equations and permits the single gauge condition.
\end{proof}

For a Latin cell variable $X_{rcs}$ asserting $L(r,c)=s$, the
antecedents $A,B$ for $f_a(i)=v$ and $f_b(i)=v$ are:
\begin{center}
\begin{tabular}{ccccc}
\toprule
View & Position & Value & $A$ & $B$\\
\midrule
row & column $i$ & symbol $v$ & $X_{aiv}$ & $X_{biv}$\\
column & row $i$ & symbol $v$ & $X_{iav}$ & $X_{ibv}$\\
symbol & column $i$ & row $v$ & $X_{via}$ & $X_{vib}$\\
\bottomrule
\end{tabular}
\end{center}
Writing $e=e_i$ and $h=h_v$, four clauses implement the equations:
\[
 (\neg A\vee\neg e\vee h),\quad(\neg A\vee e\vee\neg h),\quad
 (\neg B\vee e\vee h),\quad(\neg B\vee\neg e\vee\neg h).
\]
With Latin exactly-one constraints, imposing these clauses
independently for every line pair in every view is equivalent to FFF.
The reverse implication extends each FFF table by its pairwise
alternating colorings; no coupling of arbitrary gauges is required.
Isotopy invariance permits reduced tables in an existence search.

Each line pair uses $2n$ colors, $4n^2$ conditional clauses, and
one optional gauge unit. Balance follows from the equations and
requires no cardinality auxiliaries. Frozen positive controls at
orders $8$, $12$, and $14$ and checked negative order-$6$ controls
validate the implementation. In one fixed order-$18$ row context,
the recorded model has $22356$ variables and $745065$ clauses,
compared with $270216$ and $3190617$ in the prior physical-arc
representation. Both bounded $60$-second solver probes timed out.
This is a size reduction, not evidence of faster solution or of
nonexistence at order $18$.
